% Source: https://artofproblemsolving.com/community/c3h2297889_proving_theorems
% Stable identity: aops-topic-2297889-q01
Let $\angle1$ and $\angle2$ be right angles.

\begin{asy}
usepackage("amsmath");
size(4.5cm);
defaultpen(fontsize(10pt));
pair B=(0,0); draw((1.6,0)--B--(0,1.4),linewidth(0.9)); draw((0,0.25)--(0.25,0.25)--(0.25,0));
label("$\angle1$",(0.45,0.45),NE,fontsize(9pt));
pair E0=(3,0); draw(E0+1.6*dir(35)--E0--E0+1.4*dir(125),linewidth(0.9));
draw(E0+0.25*dir(35)--E0+0.25*dir(35)+0.25*dir(125)--E0+0.25*dir(125));
label("$\angle2$",E0+0.6*dir(80),N,fontsize(9pt));
\end{asy}

\begin{tabular}{ll}
\textbf{Statements}&\textbf{Reasons}\\\hline
$\angle1$ and $\angle2$ are right angles&Given\\
$m\angle1=90^\circ$ and $m\angle2=90^\circ$&Definition of a right angle\\
$m\angle1=m\angle2$&Transitive property of equality\\
$\angle1\cong\angle2$&Definition of congruent angles
\end{tabular}

So any two right angles are congruent. $\blacksquare$
